\chapter{快速编译与测试}

\section{下载}

下载软件包，并切换到自己的Git分支：

\begin{lstlisting}[language=bash]
git clone https://github.com/HanHFEM/mumps_cxx.git
cd mumps_cxx
git checkout -b xxx
\end{lstlisting}

如果要把本地分支上传到远端，指定远端和分支：

\begin{lstlisting}[language=bash]
git push origin xxx
\end{lstlisting}

\section{编译}

\subsection{安装必需库}

安装Fortran编译器：

\begin{lstlisting}[language=bash]
sudo apt-get install gfortran
\end{lstlisting}

安装OpenMPI：

\begin{lstlisting}[language=bash]
sudo apt-get install libopenmpi-dev
\end{lstlisting}

安装ScaLAPACK、METIS、ParMETIS和LAPACK：

\begin{lstlisting}[language=bash]
sudo apt-get install libscalapack-openmpi-dev
sudo apt-get install libscalapack-mpi-dev
sudo apt-get install libmetis-dev
sudo apt-get install libparmetis-dev
sudo apt-get install liblapack-dev
\end{lstlisting}

不同Ubuntu版本提供的ScaLAPACK包可能不同。若两个ScaLAPACK包互相冲突，应根据当前MPI实现选择其中一个。


\subsection{配置附加选项}

进入MUMPS目录：

\begin{lstlisting}[language=bash]
cd mumps-5.6.2.2
\end{lstlisting}

可以参考项目中的以下说明文件：

\begin{itemize}
  \item \texttt{Readme\_options.md}；
  \item \texttt{Readme\_ordering.md}；
  \item \texttt{Readme\_LAPCK.md}。
\end{itemize}

默认使用的配置命令为：

\begin{lstlisting}[language=bash]
cmake -B build \
  -Dparmetis=yes \
  -DCMAKE_INSTALL_PREFIX=./install \
  -DCMAKE_PREFIX_PATH=/mingw64 \
  -Dparallel=true
\end{lstlisting}

参数含义：

\begin{itemize}
  \item \texttt{-Dparmetis=yes}：启用ParMETIS等相关功能；
  \item \texttt{-DCMAKE\_INSTALL\_PREFIX=./install}：把安装目录设为当前目录下的 \texttt{install}；
  \item \texttt{-DCMAKE\_PREFIX\_PATH=/mingw64}：增加依赖库搜索前缀，该路径需要根据本机环境调整；
  \item \texttt{-Dparallel=true}：启用MPI并行版本。
\end{itemize}

\subsection{执行构建}

\begin{lstlisting}[language=bash]
cmake --build build
\end{lstlisting}

进入 \texttt{mumps-5.6.2.2/build} 后，可以看到典型构建产物：

\begin{lstlisting}
libdmumps.a
libmumps_common.a
libpord.a
libsmumps.a
以及一系列.mod文件
\end{lstlisting}

\section{编译C/C++算例}

\subsection{修改测试目录的CMakeLists.txt}

在 \texttt{mumps\_cxx/mumps-5.6.2.2/test} 中增加：

\begin{itemize}
  \item \texttt{compile\_test.c}；
  \item \texttt{cpp\_example.cpp}。
\end{itemize}

修改后的 \texttt{test/CMakeLists.txt} 整理如下：

\subsubsection{项目语言}

\begin{lstlisting}
project(MUMPS_Example LANGUAGES C CXX Fortran)
\end{lstlisting}

定义项目名称，并声明项目使用C、C++和Fortran。

\subsubsection{MPI和hwloc路径}

\begin{lstlisting}
link_directories(/usr/lib/x86_64-linux-gnu)
set(MPI_INCLUDE_DIR /usr/lib/x86_64-linux-gnu/openmpi/include)
set(MPI_C_LIBRARY /usr/lib/x86_64-linux-gnu/openmpi/lib/libmpi.so)
set(MPI_CXX_LIBRARY /usr/lib/x86_64-linux-gnu/openmpi/lib/libmpi_cxx.so)

set(HWLOC_LIBRARY /usr/lib/x86_64-linux-gnu/libhwloc.so)
set(HWLOC_INCLUDE_DIR /usr/include)
\end{lstlisting}

如果CMake因为环境或版本问题找不到某些库，可以显式指定路径。配置顶层工程时也要保证这些依赖路径可见。

\begin{lstlisting}
include_directories(${MPI_INCLUDE_DIR} ${HWLOC_INCLUDE_DIR})
\end{lstlisting}

这条命令把MPI和hwloc头文件目录加入编译器搜索路径。

\subsubsection{\texttt{mumps\_test}函数}

\begin{lstlisting}
set_property(DIRECTORY PROPERTY LABELS "unit;mumps")

function(mumps_test name tgt)
  if(parallel)
    add_test(
      NAME ${name}
      COMMAND ${MPIEXEC_EXECUTABLE} ${MPIEXEC_NUMPROC_FLAG} 2 $<TARGET_FILE:${tgt}>
    )
  else()
    add_test(NAME ${name} COMMAND ${tgt})
  endif()
endfunction()
\end{lstlisting}

\texttt{mumps\_test}接受两个参数：

\begin{itemize}
  \item \texttt{name}：CTest中的测试名称；
  \item \texttt{tgt}：需要运行的CMake可执行目标。
\end{itemize}

当 \texttt{parallel} 为真时，测试命令包含：

\begin{itemize}
  \item \texttt{\$\{MPIEXEC\_EXECUTABLE\}}：MPI启动程序；
  \item \texttt{\$\{MPIEXEC\_NUMPROC\_FLAG\}}：指定进程数的参数；
  \item \texttt{2}：启动两个MPI进程；
  \item \texttt{\$\textless{}TARGET\_FILE:\$\{tgt\}\textgreater{}}：目标可执行文件的完整路径。
\end{itemize}

当 \texttt{parallel} 为假时，CTest直接执行目标程序。

\subsubsection{双精度配置测试}

\begin{lstlisting}
if(BUILD_DOUBLE)
  add_executable(mumpscfg test_mumps.f90)
  target_link_libraries(mumpscfg PRIVATE MUMPS::MUMPS)
  mumps_test(Cfg mumpscfg)
endif()
\end{lstlisting}

当启用双精度实数版本时：

\begin{enumerate}
  \item 从 \texttt{test\_mumps.f90} 创建 \texttt{mumpscfg}；
  \item 将它链接到 \texttt{MUMPS::MUMPS}；
  \item 注册名为 \texttt{Cfg} 的测试。
\end{enumerate}

\subsubsection{版本检查}

\begin{lstlisting}
if(MUMPS_UPSTREAM_VERSION VERSION_LESS 5.1)
  return()
endif()
\end{lstlisting}

如果上游MUMPS版本低于5.1，则停止处理当前CMakeLists。

\subsubsection{单精度和双精度Fortran算例}

\begin{lstlisting}
if(BUILD_SINGLE)
  add_executable(s_simple s_simple.f90)
  target_link_libraries(s_simple PRIVATE MUMPS::MUMPS)
  mumps_test(SimpleReal32 s_simple)
endif()

if(BUILD_DOUBLE)
  add_executable(d_simple d_simple.f90)
  target_link_libraries(d_simple PRIVATE MUMPS::MUMPS)
  mumps_test(SimpleReal64 d_simple)
endif()
\end{lstlisting}

\begin{itemize}
  \item \texttt{s\_simple}对应单精度实数版本；
  \item \texttt{d\_simple}对应双精度实数版本。
\end{itemize}

\subsubsection{C算例}

\begin{lstlisting}
add_executable(Csimple simple.c)
target_link_libraries(Csimple PRIVATE MUMPS::MUMPS)
mumps_test(CsimpleReal64 Csimple)
\end{lstlisting}

这三行从 \texttt{simple.c}创建 \texttt{Csimple}，链接MUMPS，并注册一个双精度实数C接口测试。

\subsubsection{新增C/C++算例}

\begin{lstlisting}
add_executable(compile_test compile_test.c)
target_link_libraries(compile_test PRIVATE MUMPS::MUMPS)

add_executable(cpp_example cpp_example.cpp)
target_link_libraries(cpp_example PRIVATE
  ${HWLOC_LIBRARY}
  ${MPI_C_LIBRARY}
  ${MPI_CXX_LIBRARY}
  MUMPS::MUMPS
)
\end{lstlisting}

C++程序显式链接了hwloc、MPI C库、MPI C++库和MUMPS。

\subsubsection{测试属性和Windows DLL}

\begin{lstlisting}
get_property(test_names DIRECTORY ${CMAKE_CURRENT_SOURCE_DIR} PROPERTY TESTS)
set_property(TEST ${test_names} PROPERTY RESOURCE_LOCK cpu_mpi)
set_property(TEST ${test_names} PROPERTY TIMEOUT 30)
\end{lstlisting}

\begin{itemize}
  \item \texttt{get\_property}取得当前目录注册的全部测试；
  \item \texttt{RESOURCE\_LOCK cpu\_mpi}防止多个测试同时占用同一MPI测试资源；
  \item \texttt{TIMEOUT 30}把超时时间设为30秒。
\end{itemize}

\begin{lstlisting}
if(WIN32 AND BUILD_SHARED_LIBS)
  set_property(TEST ${test_names} PROPERTY
    ENVIRONMENT_MODIFICATION
    "PATH=path_list_append:${CMAKE_INSTALL_FULL_BINDIR};PATH=path_list_append:${CMAKE_PREFIX_PATH}/bin;PATH=path_list_append:${PROJECT_BINARY_DIR}"
  )
endif()
\end{lstlisting}

这段只在Windows共享库构建中生效，用于把DLL所在目录加入测试进程的 \texttt{PATH}。

构建完成后，需要在 \texttt{build}目录中找到生成的可执行文件，或者使用 \texttt{ctest}运行已注册测试。

\subsection{新建一个测试目录}

\subsubsection{继续使用CTest框架}

在 \texttt{mumps\_cxx}下新建目录，放入一个C源文件和一个简化的 \texttt{CMakeLists.txt}：

\begin{lstlisting}
set_property(DIRECTORY PROPERTY LABELS "unit;mumps")

function(mumps_test name tgt)
  if(parallel)
    add_test(
      NAME ${name}
      COMMAND ${MPIEXEC_EXECUTABLE} ${MPIEXEC_NUMPROC_FLAG} 2 $<TARGET_FILE:${tgt}>
    )
  else()
    add_test(NAME ${name} COMMAND ${tgt})
  endif()
endfunction()

if(MUMPS_UPSTREAM_VERSION VERSION_LESS 5.1)
  return()
endif()

if(BUILD_DOUBLE)
  add_executable(dotest testano.c)
  target_link_libraries(dotest PRIVATE MUMPS::MUMPS)
  mumps_test(doTEST dotest)
endif()

get_property(test_names DIRECTORY ${CMAKE_CURRENT_SOURCE_DIR} PROPERTY TESTS)
set_property(TEST ${test_names} PROPERTY RESOURCE_LOCK cpu_mpi)
set_property(TEST ${test_names} PROPERTY TIMEOUT 30)
\end{lstlisting}

还要在顶层 \texttt{CMakeLists.txt}中加入该子目录：

\begin{lstlisting}
include(cmake/summary.cmake)
include(cmake/install.cmake)

add_subdirectory(coderun)

file(GENERATE OUTPUT .gitignore CONTENT "*")
\end{lstlisting}

重新构建后，可使用 \texttt{ctest}运行新程序。

\subsubsection{构建完成后直接运行}

如果不想使用CTest，可以为目标添加构建后命令：

\begin{lstlisting}
set_property(DIRECTORY PROPERTY LABELS "unit;mumps")

if(MUMPS_UPSTREAM_VERSION VERSION_LESS 5.1)
  return()
endif()

if(BUILD_DOUBLE)
  add_executable(dotest testano.c)
  target_link_libraries(dotest PRIVATE MUMPS::MUMPS)

  add_custom_command(
    TARGET dotest POST_BUILD
    COMMAND ${CMAKE_COMMAND} -E echo "Running dotest..."
    COMMAND $<TARGET_FILE:dotest>
    DEPENDS dotest
  )
endif()
\end{lstlisting}

这样执行 \texttt{cmake --build build}时，不仅会构建 \texttt{dotest}，还会在构建完成后立即运行它。

\subsubsection{数值类型构建选项}

\texttt{mumps-5.6.2.2/options.cmake}中包含：

\begin{lstlisting}
option(BUILD_SINGLE "Build single precision float32 real" ON)
option(BUILD_DOUBLE "Build double precision float64 real" ON)
option(BUILD_COMPLEX "Build single precision complex")
option(BUILD_COMPLEX16 "Build double precision complex")
\end{lstlisting}

这些选项分别控制：

\begin{longtable}{p{0.30\textwidth}p{0.62\textwidth}}
\toprule
\textbf{选项} & \textbf{数值类型} \\
\midrule
\endfirsthead
\toprule
\textbf{选项} & \textbf{数值类型} \\
\midrule
\endhead
\texttt{BUILD\_SINGLE} & 单精度实数 \\
\texttt{BUILD\_DOUBLE} & 双精度实数 \\
\texttt{BUILD\_COMPLEX} & 单精度复数 \\
\texttt{BUILD\_COMPLEX16} & 双精度复数 \\
\bottomrule
\end{longtable}

配置中启用了单精度和双精度实数版本，没有启用复数版本。

\section{算例编写}

\subsection{\texttt{simple.c}研究}

\texttt{test/simple.c}使用DMUMPS的C接口求解：

\[
A=\begin{bmatrix}1&0\\0&2\end{bmatrix},\qquad
b=\begin{bmatrix}1\\4\end{bmatrix},\qquad
x=\begin{bmatrix}1\\2\end{bmatrix}.
\]

完整代码如下：

\subsubsection{头文件}

\begin{lstlisting}[language=C]
#include <stdio.h>
#include <string.h>
#include "mpi.h"
#include "dmumps_c.h"
\end{lstlisting}

\texttt{dmumps\_c.h}位于：

\begin{lstlisting}
mumps-5.6.2.2/build/_deps/mumps-src/include/dmumps_c.h
\end{lstlisting}

它定义了双精度实数MUMPS的C接口，包括：

\begin{itemize}
  \item \texttt{DMUMPS\_STRUC\_C}参数结构体；
  \item \texttt{dmumps\_c()}入口函数。
\end{itemize}

\begin{lstlisting}[language=C]
void MUMPS_CALL dmumps_c(DMUMPS_STRUC_C *dmumps_par);
\end{lstlisting}

\subsubsection{宏常量}

\begin{lstlisting}[language=C]
#define JOB_INIT -1
#define JOB_END -2
#define USE_COMM_WORLD -987654
\end{lstlisting}

\begin{itemize}
  \item \texttt{JOB\_INIT}：初始化MUMPS实例；
  \item \texttt{JOB\_END}：终止MUMPS实例；
  \item \texttt{USE\_COMM\_WORLD}：使用 \texttt{MPI\_COMM\_WORLD}。
\end{itemize}

\subsubsection{矩阵和右端项}

\begin{lstlisting}[language=C]
DMUMPS_STRUC_C id;
MUMPS_INT n = 2;
MUMPS_INT8 nnz = 2;
MUMPS_INT irn[] = {1, 2};
MUMPS_INT jcn[] = {1, 2};
double a[2];
double rhs[2];
\end{lstlisting}

\begin{itemize}
  \item \texttt{n}：矩阵阶数；
  \item \texttt{nnz}：非零元数量；
  \item \texttt{irn/jcn}：非零元的1-based全局行列索引；
  \item \texttt{a}：非零元数值；
  \item \texttt{rhs}：右端项，求解完成后集中式解也写回该数组。
\end{itemize}

\texttt{MUMPS\_INT8}表示MUMPS使用的64位整数类型，不是8位整数。启用 \texttt{-DINTSIZE64}时，\texttt{MUMPS\_INT}也会变成64位；某些MPI C接口参数仍要求标准C \texttt{int}。

\subsubsection{初始化MPI和MUMPS}

\begin{lstlisting}[language=C]
ierr = MPI_Init(&argc, &argv);
ierr = MPI_Comm_rank(MPI_COMM_WORLD, &myid);

id.comm_fortran = USE_COMM_WORLD;
id.par = 1;
id.sym = 0;
id.job = JOB_INIT;
dmumps_c(&id);
\end{lstlisting}

\begin{itemize}
  \item \texttt{MPI\_Init}初始化MPI环境；
  \item \texttt{MPI\_Comm\_rank}取得当前rank；
  \item \texttt{comm\_fortran}指定MUMPS通信器；
  \item \texttt{par=1}表示host也参与分解和求解；
  \item \texttt{sym=0}选择非对称路径。虽然示例矩阵实际为对称对角阵，选择非对称路径仍可求解；
  \item \texttt{job=-1}初始化MUMPS实例。
\end{itemize}

\subsubsection{只在host上提供集中式矩阵}

\begin{lstlisting}[language=C]
if (myid == 0) {
  id.n = n;
  id.nnz = nnz;
  id.irn = irn;
  id.jcn = jcn;
  id.a = a;
  id.rhs = rhs;
}
\end{lstlisting}

默认集中式输入模式下，rank 0提供矩阵和右端项，其他rank参与MUMPS内部并行计算。

\subsubsection{C数组下标与手册编号}

\begin{lstlisting}[language=C]
#define ICNTL(I) icntl[(I)-1]
\end{lstlisting}

手册中的控制参数从1开始编号，而C数组从0开始。该宏使代码可以继续使用手册中的编号，例如 \texttt{id.ICNTL(1)}实际访问 \texttt{id.icntl[0]}。

\begin{lstlisting}[language=C]
id.ICNTL(1) = -1;
id.ICNTL(2) = -1;
id.ICNTL(3) = -1;
id.ICNTL(4) = 0;
\end{lstlisting}

示例关闭了主要输出信息。

\subsubsection{分析、分解和求解}

\begin{lstlisting}[language=C]
id.job = 6;
dmumps_c(&id);
\end{lstlisting}

\texttt{JOB=6}等价于依次执行：

\begin{enumerate}
  \item 分析；
  \item 数值分解；
  \item 求解。
\end{enumerate}

\subsubsection{错误检查和释放}

\begin{lstlisting}[language=C]
if (id.infog[0] < 0) {
  fprintf(stderr,
    " (PROC %d) ERROR RETURN: INFOG(1)= %d INFOG(2)= %d\n",
    myid, id.infog[0], id.infog[1]);
  return 1;
}

id.job = JOB_END;
dmumps_c(&id);
MPI_Finalize();
\end{lstlisting}

\texttt{INFOG(1)\textless{}0}表示调用失败，\texttt{INFOG(2)}通常提供附加错误信息。程序结束前先以 \texttt{JOB=-2}释放MUMPS实例，再调用 \texttt{MPI\_Finalize}。

\subsection{JOB参数}

在调用MUMPS前，所有MPI进程都要设置相同的 \texttt{JOB}。MUMPS不会修改该字段。

\begin{longtable}{p{0.30\textwidth}p{0.62\textwidth}}
\toprule
\textbf{\texttt{JOB}} & \textbf{作用} \\
\midrule
\endfirsthead
\toprule
\textbf{\texttt{JOB}} & \textbf{作用} \\
\midrule
\endhead
\texttt{-1} & 初始化；调用前先设置 \texttt{comm\_fortran}、\texttt{sym}和 \texttt{par} \\
\texttt{-2} & 终止实例并释放内部资源 \\
\texttt{1} & 分析矩阵结构 \\
\texttt{2} & 数值分解，要求已有成功的分析结果 \\
\texttt{3} & 求解，要求已有成功的数值分解 \\
\texttt{4} & 分析和数值分解 \\
\texttt{5} & 数值分解和求解 \\
\texttt{6} & 分析、数值分解和求解 \\
\texttt{7} & 把MUMPS内部实例保存到磁盘 \\
\texttt{-3} & 删除保存到磁盘的实例数据 \\
\texttt{8} & 从磁盘恢复MUMPS实例 \\
\texttt{-4} & 释放数值阶段数据，但保留分析阶段结构 \\
\texttt{9} & 在求解前计算分布式右端项的可能分布 \\
\bottomrule
\end{longtable}

\texttt{JOB=9}只能在成功完成分解，即 \texttt{JOB=2/4}返回且 \texttt{INFOG(1)\textgreater{}=0}后使用。

\subsection{PAR参数}

\texttt{PAR}的取值为：

\begin{itemize}
  \item \texttt{PAR=0}：host不参加分解和求解的并行计算，主要负责初始问题、分析协调、数据分发和结果收集；
  \item \texttt{PAR=1}：host也参加分解和求解；其他值按1处理。
\end{itemize}

所有进程应提供相同的 \texttt{PAR}。如果不一致，MUMPS使用rank 0上的值。

当矩阵集中在rank 0且输入规模很大时，\texttt{PAR=1}可能让rank 0同时承担原始矩阵和数值工作区，从而造成内存不均衡。

\subsection{矩阵输入格式}

详细接口可以参考 MUMPS Users' Guide 第5.2节。

\subsubsection{三类输入结构}

\begin{lstlisting}[language=C]
/* Assembled entry */
MUMPS_INT      nz;
MUMPS_INT8     nnz;
MUMPS_INT      *irn;
MUMPS_INT      *jcn;
DMUMPS_COMPLEX *a;

/* Distributed entry */
MUMPS_INT      nz_loc;
MUMPS_INT8     nnz_loc;
MUMPS_INT      *irn_loc;
MUMPS_INT      *jcn_loc;
DMUMPS_COMPLEX *a_loc;

/* Element entry */
MUMPS_INT      nelt;
MUMPS_INT      *eltptr;
MUMPS_INT      *eltvar;
DMUMPS_COMPLEX *a_elt;
\end{lstlisting}

\subsubsection{SYM}

\texttt{SYM}指定矩阵类型：

\begin{itemize}
  \item \texttt{0}：非对称矩阵；
  \item \texttt{1}：对称正定矩阵，使用不进行数值主元搜索的正定路径；
  \item \texttt{2}：一般对称矩阵；
  \item 其他值按0处理。
\end{itemize}

所有进程应提供相同的 \texttt{SYM}，否则使用rank 0上的值。MUMPS不提供专门的Hermitian矩阵模式。
